%% kapa: exam=1 round=29 year=2018 period=2 form=A subject=appraisal_law subject_ko=감정평가관계법규 seq=13 number=13
%% source: Q-Net 감정평가사 기출문제 (cst00309) artlSeq=5202808 file=extracted/1008/5202808/2018년 제29회 감정평가사 1차시험 기출문제/2018년 제29회 감정평가사 1차시험 기출문제/2교시 A형.pdf
%% conversion: kordoc 4.24.51 -> tools/gen1.py (auto); review=proofread-sonnet; flags=none
\begin{question}{13}
국토의 계획 및 이용에 관한 법령상 도시ㆍ군계획시설사업에 관한 설명으로 옳지 않은 것은?
\begin{choices}
  \item 도시ㆍ군계획시설사업의 시행자는 사업시행대상지역 또는 대상시설을 둘 이상으로 분할하여 도시ㆍ군계획시설사업을 시행할 수 있다.
  \item 한국토지주택공사가 도시ㆍ군계획시설사업시행자로 지정을 받으려면 토지소유자 총수의 2분의 1 이상에 해당하는 자의 동의를 얻어야 한다.
  \item 도지사는 광역도시계획과 관련되거나 특히 필요하다고 인정되는 경우에는 관계 시장 또는 군수의 의견을 들어 직접 도시ㆍ군계획시설사업을 시행할 수 있다.
  \item 도시ㆍ군계획시설에 대한 단계별 집행계획은 제1단계 집행계획과 제2단계 집행계획으로 구분하여 수립하되, 3년 이내에 시행하는 도시ㆍ군계획시설사업은 제1단계 집행계획에 포함되도록 하여야 한다.
  \item 도시ㆍ군계획시설사업의 시행자로 지정받은「지방공기업법」에 의한 지방공사는 기반시설의 설치가 필요한 경우에 그 이행을 담보하기 위한 이행보증금을 예치하지 않아도 된다.
\end{choices}
\end{question}
\answer{2}
